\begin{figure}[ht]
\centering
\begin{tikzpicture}[x=1cm, y=1cm, font=\small]

% Convergence comparison: log10 of the absolute error versus iteration
% index for three root-finders on a smooth root. Newton (quadratic),
% secant (order ~1.618), bisection (linear, slope ~ -log10(2) per step).
% Axes: iteration 0..8 mapped to x in [0, 9.6] (1.2 per step);
% log10 error from 0 down to -16 mapped to y: y = 0 at top (log err 0),
% y = 4.8 at bottom is log err -16, so y = 4.8 + 0.3*logerr.

% Frame
\draw[->, thick] (-0.3, 0) -- (10.2, 0) node[right] {iteration $n$};
\draw[->, thick] (0, -0.3) -- (0, 5.4)
  node[above, align=center] {$\log_{10}|e_{n}|$};

% y ticks (log error values 0, -4, -8, -12, -16)
\foreach \le in {0, -4, -8, -12, -16} {
  \pgfmathsetmacro\yy{4.8 + 0.3*\le}
  \draw (0, \yy) -- (-0.12, \yy);
  \node[anchor=east, font=\scriptsize] at (-0.14, \yy) {\le};
}
% x ticks
\foreach \n in {0,1,2,3,4,5,6,7,8} {
  \pgfmathsetmacro\xx{1.2*\n}
  \draw (\xx, 0) -- (\xx, -0.12);
  \node[anchor=north, font=\scriptsize] at (\xx, -0.12) {\n};
}

% Newton: log err doubles in magnitude each step (quadratic).
% logerr_n approx -1 * 2^n, clipped at -16.
% n: 0->-0.5, 1->-1, 2->-2, 3->-4, 4->-8, 5->-16
\draw[very thick, engblue]
  plot coordinates {
    ({1.2*0}, {4.8 + 0.3*(-0.5)})
    ({1.2*1}, {4.8 + 0.3*(-1)})
    ({1.2*2}, {4.8 + 0.3*(-2)})
    ({1.2*3}, {4.8 + 0.3*(-4)})
    ({1.2*4}, {4.8 + 0.3*(-8)})
    ({1.2*5}, {4.8 + 0.3*(-16)})
  };
\foreach \pt in {{0/-0.5},{1/-1},{2/-2},{3/-4},{4/-8},{5/-16}} {}
\node[engblue, anchor=west, font=\scriptsize] at (6.4, 4.6)
  {Newton (order 2)};

% Secant: order phi approx 1.618. logerr_n approx -0.5 * 1.618^n.
% n: 0->-0.5, 1->-0.81, 2->-1.31, 3->-2.12, 4->-3.43, 5->-5.55,
% 6->-8.98, 7->-14.5, 8->-16(clip)
\draw[very thick, engorange, dashed]
  plot coordinates {
    ({1.2*0}, {4.8 + 0.3*(-0.5)})
    ({1.2*1}, {4.8 + 0.3*(-0.81)})
    ({1.2*2}, {4.8 + 0.3*(-1.31)})
    ({1.2*3}, {4.8 + 0.3*(-2.12)})
    ({1.2*4}, {4.8 + 0.3*(-3.43)})
    ({1.2*5}, {4.8 + 0.3*(-5.55)})
    ({1.2*6}, {4.8 + 0.3*(-8.98)})
    ({1.2*7}, {4.8 + 0.3*(-14.5)})
    ({1.2*8}, {4.8 + 0.3*(-16)})
  };
\node[engorange, anchor=west, font=\scriptsize] at (6.4, 4.2)
  {secant (order $\varphi$)};

% Bisection: linear, logerr decreases by log10(2)=0.301 per step.
% n: 0->-0.5, then -0.5 - 0.301*n
\draw[very thick, engred, dotted]
  plot coordinates {
    ({1.2*0}, {4.8 + 0.3*(-0.5)})
    ({1.2*1}, {4.8 + 0.3*(-0.80)})
    ({1.2*2}, {4.8 + 0.3*(-1.10)})
    ({1.2*3}, {4.8 + 0.3*(-1.40)})
    ({1.2*4}, {4.8 + 0.3*(-1.70)})
    ({1.2*5}, {4.8 + 0.3*(-2.00)})
    ({1.2*6}, {4.8 + 0.3*(-2.31)})
    ({1.2*7}, {4.8 + 0.3*(-2.61)})
    ({1.2*8}, {4.8 + 0.3*(-2.91)})
  };
\node[engred, anchor=west, font=\scriptsize] at (6.4, 3.8)
  {bisection (linear)};

\end{tikzpicture}
\caption[Convergence rates of Newton, secant, and bisection]{Absolute
error against iteration index on a smooth root, plotted on a
logarithmic error axis. Newton's method doubles the number of correct
digits per step: the error curve bends downward, each segment twice as
steep as the last, reaching machine precision in about five steps from
a starting error of $10^{-0.5}$. The secant method, which replaces the
analytic derivative with a finite difference across the last two
iterates, converges with order $\varphi = (1 + \sqrt{5})/2 \approx
1.618$, slightly slower per step but free of a derivative evaluation.
Bisection halves the error every step, a straight line of slope
$-\log_{10} 2 \approx -0.30$, and reaches the same precision only after
roughly fifty steps. The superlinear methods win decisively once the
iterate is inside the basin where their error models hold.}
\label{fig:vol02:ch05:secant-convergence}
\end{figure}
